% Paragraphs 1-3 are Dr. Lin's supplied text, installed VERBATIM on 2026-09-08. Only the
% bracketed reference numbers are mapped to bibliography keys; not one word is changed.
% Grammar, factual and citation concerns are recorded in
% pipeline/reports/POST_MEETING_REVIEW_HANDOFF.md, not corrected here.
% Paragraph 4 and the contributions are the framework/contribution continuation the meeting
% asked for; they are not protected text.
\section{Introduction}
\label{sec:intro}

Device fingerprinting has been an essential step in cyber reconnaissance, allowing adversaries
to reveal unique features of target networks and to design effective, stealthy attack
strategies. These techniques are becoming increasingly critical in industrial control systems
(ICSs) such as power grids, where adversaries often use IP-based control networks and computing
devices within as entry points to inflict physical damage. Consequently, fingerprinting shifts
the focus from visited websites and user biometric behavior to device models and types of
control operations that are critical to ICS attacks. In the 2015 attack that disrupted Ukrainian
power grids and the Stuxnet attack that disrupted Iranian nuclear power facilities, it is widely
believed that adversaries stayed in their systems for at least 6 months to perform cyber
reconnaissance~\cite{ukraine2015, stuxnet2011}.

To disrupt device fingerprinting, many studies have presented network traffic
obfuscation~\cite{wright2009, dyer2012}. Device
fingerprinting targeting general computing environments generally relies on network-level
features such as packet size and/or inter-packet latency observed from communication
patterns~\cite{sirinam2018}. Consequently, existing traffic
obfuscation often focuses on (i) padding and splitting network packets, which hide or change the
distribution of network packet sizes~\cite{wright2009, walkietalkie2017, panchenko2016}, and (ii)
delaying network packets and adding dummy ones, which disrupt the inter-packet timing
pattern~\cite{dyer2012, tamaraw2014, wang2008dlp, juarez2016}. Since manipulating
communication networks can introduce runtime overhead, recent works have begun to offload
traffic obfuscation onto programmable network switches, which change communication patterns at
much higher line rates than CPUs~\cite{ditto2022, minos2025, securitas2026}.

Unfortunately, these methods cannot be directly applied to ICS environments for two main
reasons. First, device fingerprinting methods on ICS devices rely on different features from the
ones in general computing environments. Instead of using network layer knowledge, they attempt
to use behavior on physical devices to reveal device model or types of operations. For example,
work in~\cite{formby2016, shen2018, ahmed2021, pospisil2023} uses the latency between the
TCP acknowledge packets and the actual response from the same device to estimate execution time
there. Since each
vendors may choose different materials or algorithms to implement the same control
functionality, this time can be accurate in fingerprinting those devices. Second, existing
traffic obfuscation is performed over encrypted channels, which are necessary to mix the
obfuscated and normal traffic. Unfortunately, ICS networks still rely on unencrypted traffic due
to a wide range of legacy devices, despite ICS network protocols may provide security features.
Directly padding or adding dummy packets can easily reveal obfuscated packets, exposing the the
underlying traffic pattern to adversaries.

In this paper, we develop an operation-aware framework for timing and size
policies on a programmable switch. The framework identifies protocol control
units and applies supported padding, splitting and scheduling policies while
retaining legitimate endpoint communication. Case 4 is the new implementation
target. Our current evidence establishes bounded software composition and a
compiled recovery component; the complete joint Tofino realization remains an
implementation gate. These claims require neither endpoint encryption nor
separate de-obfuscation hardware, and do not establish hidden device identity.
